\chapter{Academic Documents}
\label{ch:academic}

This chapter covers five doctypes designed for teaching and learning
contexts: \texttt{homework}, \texttt{exam}, \texttt{lecture-notes},
\texttt{syllabus}, and \texttt{handout}.  All use \texttt{scrartcl} with
serif fonts, one-and-a-half line spacing, and half paragraph spacing.

\section{Homework}
\label{sec:homework}

The \texttt{homework} doctype produces a title page with course name,
assignment number, student details, due date, and total points.
It suppresses the TOC and sets \texttt{secnumdepth=0} for a clean,
assignment-style layout.

\begin{minted}{latex}
\documentclass[doctype=homework]{../../omnilatex}
\title{Linear Algebra Problem Set}
\hwcourse{MATH 201}
\hwnumber{Homework 3}
\hwauthor{Prof.\ A.\ Smith}
\hwduedate{2025-11-15}
\hwinstructor{Prof.\ A.\ Smith}
\hwstudentname{Jane Doe}
\hwstudentid{12345678}
\hwtotalpoints{50}
\begin{document}
\maketitle
\begin{exercise}[10]
    Prove that every vector space has a basis.
\end{exercise}
\begin{solution}
    By Zorn's Lemma, \ldots
\end{solution}
\end{document}
\end{minted}

\subsection{Custom Commands}

\begin{table}[htbp]
    \centering
    \caption{Homework metadata commands}
    \label{tab:homework-commands}
    \begin{tabular}{@{} l p{7.5cm} @{}}
        \toprule
        \textbf{Command} & \textbf{Description} \\
        \midrule
        \texttt{\textbackslash hwtitle\{...\}}       & Assignment title (defaults to \texttt{\textbackslash title}) \\
        \texttt{\textbackslash hwcourse\{...\}}      & Course name (default: Course Name) \\
        \texttt{\textbackslash hwauthor\{...\}}      & Author name (defaults to \texttt{\textbackslash author}) \\
        \texttt{\textbackslash hwduedate\{...\}}     & Due date (defaults to \texttt{\textbackslash date}) \\
        \texttt{\textbackslash hwinstructor\{...\}}  & Instructor name (default: Instructor Name) \\
        \texttt{\textbackslash hwnumber\{...\}}      & Assignment label (default: Homework 1) \\
        \texttt{\textbackslash hwstudentname\{...\}} & Student name (default: empty) \\
        \texttt{\textbackslash hwstudentid\{...\}}   & Student ID (default: empty) \\
        \texttt{\textbackslash hwtotalpoints\{...\}} & Total points (default: empty) \\
        \bottomrule
    \end{tabular}
\end{table}

\subsection{Environments}

Two environments are provided for structuring problem sets:

\begin{itemize}
    \item \texttt{exercise[points]} --- Numbers exercises automatically and
          optionally displays a point value in parentheses.
    \item \texttt{solution} --- Wraps the solution text.  Content is only
          visible when solutions are enabled (see below).
\end{itemize}

\subsection{Solution Control}

A boolean flag controls whether solutions are rendered:

\begin{minted}{latex}
\showsolutions   % render all solution environments
\hidesolutions   % hide all solution environments (default)
\end{minted}

When \texttt{\textbackslash hidesolutions} is active (the default), content
inside \texttt{solution} environments is silently discarded via
\texttt{\textbackslash setbox}, producing a clean student-facing handout.

\section{Examinations}
\label{sec:exam}

The \texttt{exam} doctype formats formal examination papers with a title
page showing exam code, date, duration, total marks, department, and
custom instructions.  It uses \texttt{pagestyle=empty} and suppresses
section numbering.

\begin{minted}{latex}
\documentclass[doctype=exam]{../../omnilatex}
\title{Final Examination}
\examcode{MATH301-F2025}
\examdate{2025-12-18}
\examduration{3 hours}
\totalmarks{100}
\department{Department of Mathematics}
\begin{document}
\maketitle
\begin{question}[25]
    State and prove the Fundamental Theorem of Calculus.
\end{question}
\begin{answer}[4cm]
    % blank space for student response
\end{answer}
\end{document}
\end{minted}

\begin{table}[htbp]
    \centering
    \caption{Examination commands}
    \label{tab:exam-commands}
    \begin{tabular}{@{} l p{7.5cm} @{}}
        \toprule
        \textbf{Command} & \textbf{Description} \\
        \midrule
        \texttt{\textbackslash examtitle\{...\}}        & Exam title (defaults to \texttt{\textbackslash title}) \\
        \texttt{\textbackslash examcode\{...\}}         & Exam code (default: EXAM000) \\
        \texttt{\textbackslash examdate\{...\}}         & Exam date (defaults to \texttt{\textbackslash date}) \\
        \texttt{\textbackslash examduration\{...\}}     & Duration (default: 2 hours) \\
        \texttt{\textbackslash totalmarks\{...\}}       & Total marks (default: 100) \\
        \texttt{\textbackslash department\{...\}}       & Department name (default: Department of Mathematics) \\
        \texttt{\textbackslash examinstructions\{...\}} & Custom instructions \\
        \bottomrule
    \end{tabular}
\end{table}

Two environments structure exam content:
\texttt{question[marks]} auto-numbers questions and displays an optional
mark allocation, while \texttt{answer[height]} creates an
\texttt{addmargin} of configurable height for student responses.

\section{Lecture Notes}
\label{sec:lecture-notes}

The \texttt{lecture-notes} doctype is designed for textbook-style handouts
with a wide left margin for student annotations.  It sets
\texttt{secnumdepth=2}, includes a header with course name and lecture
number, and provides theorem-like environments.

\subsection{Metadata Commands}

\begin{table}[htbp]
    \centering
    \caption{Lecture notes commands}
    \label{tab:lecture-commands}
    \begin{tabular}{@{} l p{7.5cm} @{}}
        \toprule
        \textbf{Command} & \textbf{Description} \\
        \midrule
        \texttt{\textbackslash lecturetitle\{...\}} & Lecture title (defaults to \texttt{\textbackslash title}) \\
        \texttt{\textbackslash lecturenumber\{...\}} & Lecture number (default: 1) \\
        \texttt{\textbackslash coursename\{...\}} & Course name (default: Course Name) \\
        \texttt{\textbackslash lecturedate\{...\}} & Date (defaults to \texttt{\textbackslash date}) \\
        \texttt{\textbackslash lecturer\{...\}} & Lecturer name (default: Lecturer Name) \\
        \texttt{\textbackslash semester\{...\}} & Semester (default: Fall 2025) \\
        \bottomrule
    \end{tabular}
\end{table}

\subsection{Theorem-Like Environments}

Numbered environments share a common counter:
\texttt{theorem}, \texttt{lemma}, \texttt{corollary}, and
\texttt{definition}.  Unnumbered environments include \texttt{example},
\texttt{remark}, and \texttt{proof} (which appends a $\square$).

\subsection{Wide Left Margin}

The doctype calls \texttt{\textbackslash setCustomMargins\{3cm\}\{2cm\}\{2.5cm\}\{2.5cm\}},
creating a 3\,cm left margin ideal for handwritten annotations during
lectures.

\section{Syllabus}
\label{sec:syllabus}

The \texttt{syllabus} doctype generates a course syllabus with a detailed
title page listing instructor contact information, meeting details,
prerequisites, and course website.

\begin{table}[htbp]
    \centering
    \caption{Syllabus commands}
    \label{tab:syllabus-commands}
    \begin{tabular}{@{} l p{7.5cm} @{}}
        \toprule
        \textbf{Command} & \textbf{Description} \\
        \midrule
        \texttt{\textbackslash coursenumber\{...\}}    & Course code (default: CS\,000) \\
        \texttt{\textbackslash coursetitle\{...\}}     & Course title (default: Course Title) \\
        \texttt{\textbackslash semester\{...\}}        & Semester (default: Fall 2025) \\
        \texttt{\textbackslash instructor\{...\}}      & Instructor name \\
        \texttt{\textbackslash instructoroffice\{...\}} & Office location \\
        \texttt{\textbackslash instructoremail\{...\}}  & Email address \\
        \texttt{\textbackslash instructorhours\{...\}}  & Office hours \\
        \texttt{\textbackslash meetingtime\{...\}}      & Class meeting time \\
        \texttt{\textbackslash meetinglocation\{...\}}  & Classroom location \\
        \texttt{\textbackslash prerequisites\{...\}}    & Prerequisites (default: None) \\
        \texttt{\textbackslash coursewebsite\{...\}}    & Course URL \\
        \bottomrule
    \end{tabular}
\end{table}

\subsection{Environments and Helpers}

\begin{itemize}
    \item \texttt{objectives} --- Wraps a numbered list of learning objectives.
    \item \texttt{gradingpolicy} --- Produces a \texttt{booktabs} table with
          Component/Weight columns and an automatic Total row.
    \item \texttt{\textbackslash gradingitem\{component\}\{weight\}} ---
          Adds a row inside the \texttt{gradingpolicy} environment.
    \item \texttt{\textbackslash scheduleentry\{week\}\{topic\}\{notes\}} ---
          Adds a row for weekly schedule tables.
\end{itemize}

\begin{minted}{latex}
\begin{gradingpolicy}
    \gradingitem{Homework}{30}
    \gradingitem{Midterm Exam}{25}
    \gradingitem{Final Exam}{35}
    \gradingitem{Participation}{10}
\end{gradingpolicy}
\end{minted}

\section{Handout}
\label{sec:handout}

The \texttt{handout} doctype produces compact, two-column course handouts
with a running header showing course name, handout number, title, author,
and date.  It uses 15\,mm margins and \texttt{parskip=half}.

\begin{table}[htbp]
    \centering
    \caption{Handout commands}
    \label{tab:handout-commands}
    \begin{tabular}{@{} l p{7.5cm} @{}}
        \toprule
        \textbf{Command} & \textbf{Description} \\
        \midrule
        \texttt{\textbackslash handouttitle\{...\}}  & Handout title (defaults to \texttt{\textbackslash title}) \\
        \texttt{\textbackslash handoutcourse\{...\}} & Course name (default: Course Name) \\
        \texttt{\textbackslash handoutnumber\{...\}} & Handout number (default: 1) \\
        \texttt{\textbackslash handoutdate\{...\}}   & Date (defaults to \texttt{\textbackslash date}) \\
        \texttt{\textbackslash handoutauthor\{...\}} & Author (defaults to \texttt{\textbackslash author}) \\
        \bottomrule
    \end{tabular}
\end{table}

\subsection{Highlight Boxes}

Two tcolorbox-based commands draw attention to key material:

\begin{itemize}
    \item \texttt{\textbackslash keyconcept\{title\}\{text\}} --- A titled
          box with a blue frame (\texttt{cBlue!70!black}) and light
          background (\texttt{Air!40}), rendered via the \texttt{keyconceptbox}
          tcolorbox.
    \item \texttt{\textbackslash remember\{text\}} --- An inline box with an
          orange frame (\texttt{mOrange!80!black}) and yellow background
          (\texttt{mYellow!20}), prefixed with ``Remember:''.
\end{itemize}

\subsection{Two-Column Mode}

The doctype enables \texttt{twocolumn=true} via KOMA-Script, so all body
content flows in two columns automatically.  Use
\texttt{\textbackslash keyconcept} and \texttt{\textbackslash remember}
for single-column callouts that span the full column width.
